% ============================================================
%  GUÍA 3: TRANSFORMACIÓN DE UNIDADES E INTRODUCCIÓN A LOS VECTORES
%  Curso 4to Año — Física
%  Autor: Ing. Gabriel Astudillo
%  Clases cubiertas: 4 a 6
% ============================================================

\documentclass[12pt, a4paper]{article}

% ── Paquetes ─────────────────────────────────────────────────

\usepackage{mathwizards-guia}
\mwguia{Guía 3 — Unidades y Vectores}{Ing. Gabriel Astudillo $\cdot$ 4to Año}

% ── Comandos propios de esta guía ───────────────────
\providecommand{\sen}{\sin}

\begin{document}
% ============================================================

% ── Portada / Título ─────────────────────────────────────────
\mwportada{Guía de Estudio 3}{Transformación de Unidades e Introducción a los Vectores}{El Sistema Internacional, los factores de conversión y las magnitudes con dirección}

\vspace{4pt}
\begin{cuadroinfo}
  \small
  \textbf{Presentación:} ¡Hola! La física es la ciencia de \emph{medir}: medimos velocidades, masas, distancias y fuerzas. Pero medir no sirve de nada si no podemos \emph{comparar} nuestras medidas, y para eso necesitamos un idioma común de unidades. En esta guía aprenderás a moverte con soltura en el Sistema Internacional, a transformar unidades (¡como convertir 90 km/h a m/s!) y conocerás la primera gran herramienta de la física: los \textbf{vectores}, que nos permiten trabajar con magnitudes que tienen dirección, como el viento o la velocidad de un avión. ¡Allá vamos!
  \\[4pt]
  \textbf{Nivel de Dificultad:}\quad
  \facil{} Básico / Directo \quad
  \medio{} Intermedio \quad
  \dificil{} Desafiante
\end{cuadroinfo}

\vspace{8pt}

% ============================================================
\section{Sistema Internacional de Unidades y medición}
% ============================================================

\subsection{Magnitudes y unidades}

\begin{cuadroteoria}
  Una \textbf{magnitud física} es toda propiedad que se puede medir (longitud, masa, tiempo, temperatura\ldots).
  \begin{itemize}[leftmargin=*]
    \item \textbf{Magnitudes fundamentales:} no se definen a partir de otras (longitud, masa, tiempo\ldots).
    \item \textbf{Magnitudes derivadas:} se obtienen combinando fundamentales (velocidad, fuerza, densidad\ldots).
  \end{itemize}

  \textbf{Unidades fundamentales del Sistema Internacional (SI):}
  \begin{center}
    \begin{tabular}{lll}
      \toprule
      Magnitud                & Unidad    & Símbolo \\
      \midrule
      Longitud                & metro     & m       \\
      Masa                    & kilogramo & kg      \\
      Tiempo                  & segundo   & s       \\
      Intensidad de corriente & ampere    & A       \\
      Temperatura             & kelvin    & K       \\
      Cantidad de sustancia   & mol       & mol     \\
      Intensidad luminosa     & candela   & cd      \\
      \bottomrule
    \end{tabular}
  \end{center}
\end{cuadroteoria}

\begin{cuadroteoria}
  \textbf{Prefijos del SI} (multiplicadores de las unidades):
  \begin{center}
    \begin{tabular}{llll}
      \toprule
      Prefijo & Símbolo & Factor    & Ejemplo                                     \\
      \midrule
      kilo    & k       & $10^{3}$  & $1 \text{ km} = 1000 \text{ m}$             \\
      hecto   & h       & $10^{2}$  & $1 \text{ hm} = 100 \text{ m}$              \\
      deca    & da      & $10^{1}$  & $1 \text{ dam} = 10 \text{ m}$              \\
      deci    & d       & $10^{-1}$ & $1 \text{ dm} = 0{,}1 \text{ m}$            \\
      centi   & c       & $10^{-2}$ & $1 \text{ cm} = 0{,}01 \text{ m}$           \\
      mili    & m       & $10^{-3}$ & $1 \text{ mm} = 0{,}001 \text{ m}$          \\
      micro   & $\mu$   & $10^{-6}$ & $1 \text{ }\mu\text{m} = 10^{-6} \text{ m}$ \\
      \bottomrule
    \end{tabular}
  \end{center}
\end{cuadroteoria}

\begin{cuadroteoria}
  \textbf{Notación científica:} se escribe un número como $a \times 10^{n}$, con $1 \le a < 10$ y $n$ entero.
  $$4\,500\,000 = 4{,}5 \times 10^{6}, \qquad 0{,}00025 = 2{,}5 \times 10^{-4}$$

  \textbf{Cifras significativas (introducción):} son los dígitos que aportan información real de una medición. Por ejemplo, $2{,}50$ m tiene tres cifras significativas (el cero cuenta).
\end{cuadroteoria}

\begin{cuadroadvertencia}
  \textbf{Errores clásicos:}
  \begin{itemize}[leftmargin=*]
    \item Escribir $0{,}0000001$ en lugar de $10^{-7}$: la notación científica evita errores de ceros.
    \item Confundir los prefijos: $1 \text{ cm} = 0{,}01 \text{ m}$ (centi es $\frac{1}{100}$), no $0{,}1$.
    \item Escribir unidades en plural o con mayúscula incorrecta: se escribe \emph{5 m} (no "5 ms", que es milisegundo) y \emph{10 kg} (no "10 Kgs").
  \end{itemize}
\end{cuadroadvertencia}

\subsection{Ejercicios resueltos}

\textbf{Ejemplo R1.} Clasifica cada magnitud en fundamental o derivada e indica su unidad del SI: $a)$ masa \quad $b)$ velocidad \quad $c)$ tiempo.

\begin{cuadrosolucion}
  \textbf{Solución:}
  \begin{itemize}
    \item $a)$ Masa: fundamental, se mide en kilogramos (kg).
    \item $b)$ Velocidad: derivada (longitud entre tiempo), se mide en m/s.
    \item $c)$ Tiempo: fundamental, se mide en segundos (s).
  \end{itemize}
\end{cuadrosolucion}

\newpage

\textbf{Ejemplo R2.} Expresa en metros: $a)\ 5 \text{ km} \qquad b)\ 2500 \text{ mm}$.

\begin{cuadrosolucion}
  \textbf{Solución:}
  \begin{itemize}
    \item $a)\ 5 \text{ km} = 5 \times 10^{3} \text{ m} = 5000 \text{ m}$.
    \item $b)\ 2500 \text{ mm} = 2500 \times 10^{-3} \text{ m} = 2{,}5 \text{ m}$.
  \end{itemize}
\end{cuadrosolucion}

\textbf{Ejemplo R3.} Escribe en notación científica: $a)\ 0{,}00025 \qquad b)\ 4\,500\,000$.

\begin{cuadrosolucion}
  \textbf{Solución:}
  \begin{itemize}
    \item $a)\ 0{,}00025 = 2{,}5 \times 10^{-4}$ (la coma corre 4 lugares a la derecha).
    \item $b)\ 4\,500\,000 = 4{,}5 \times 10^{6}$ (la coma corre 6 lugares a la izquierda).
  \end{itemize}
\end{cuadrosolucion}

\subsection{Ejercicios propuestos}

\begin{enumerate}[leftmargin=*, resume]
  \item \facil{} Clasifica en fundamental o derivada: $a)$ longitud \quad $b)$ densidad \quad $c)$ temperatura \quad $d)$ fuerza.

  \item \facil{} Escribe la unidad del SI de: $a)$ masa \quad $b)$ tiempo \quad $c)$ longitud.

  \item \facil{} Completa: $1 \text{ km} =$ \_\_\_\_\_ m \quad ; \quad $1 \text{ cm} =$ \_\_\_\_\_ m \quad ; \quad $1 \text{ mm} =$ \_\_\_\_\_ m.

  \item \facil{} Expresa en metros: $a)\ 3 \text{ km} \qquad b)\ 750 \text{ cm} \qquad c)\ 12 \text{ mm}$.

  \item \facil{} Escribe en notación científica: $a)\ 300\,000 \qquad b)\ 0{,}0005 \qquad c)\ 25\,000\,000$.

  \item \medio{} Expresa en notación científica: $a)\ 0{,}0000002 \qquad b)\ 9\,460\,000\,000\,000$ (distancia de un año luz en km).

  \item \medio{} ¿Cuántos milímetros hay en $2{,}5$ metros? ¿Cuántos centímetros?

  \item \medio{} La velocidad del sonido es aproximadamente $340$ m/s. Exprésala en notación científica.

  \item \medio{} ¿Cuántas cifras significativas tiene cada medida? $a)\ 3{,}50 \text{ m} \qquad b)\ 0{,}0025 \text{ kg} \qquad c)\ 1200 \text{ s}$.

  \item \dificil{} Un átomo de hidrógeno mide aproximadamente $0{,}0000000001$ m. Exprésalo en notación científica y con el prefijo apropiado (angstrom $\AA = 10^{-10}$ m).

  \item \dificil{} La masa de la Tierra es $5{,}97 \times 10^{24}$ kg. Escríbela con todos sus ceros y ordénala comparándola con la masa del Sol ($1{,}99 \times 10^{30}$ kg). ¿Cuántas veces cabe la Tierra en el Sol en masa?
\end{enumerate}

\newpage

% ============================================================
\section{Transformación de unidades}
% ============================================================

\subsection{El factor de conversión}

\begin{cuadroteoria}
  Un \textbf{factor de conversión} es una fracción que vale 1 y que relaciona dos unidades equivalentes. Por ejemplo:
  $$\frac{1000 \text{ m}}{1 \text{ km}} = 1 \qquad \text{y} \qquad \frac{1 \text{ km}}{1000 \text{ m}} = 1$$
  Al multiplicar una cantidad por un factor de conversión, la cantidad no cambia: solo cambia su forma de expresarla.
\end{cuadroteoria}

\begin{cuadropasos}
  \textbf{Cómo transformar unidades:}
  \begin{enumerate}[leftmargin=*]
    \item Escribe la cantidad con su unidad.
    \item Elige el factor de conversión que \emph{cancele} la unidad de partida (la unidad de partida debe quedar en el denominador del factor).
    \item Multiplica y simplifica: los números se multiplican y las unidades se cancelan.
  \end{enumerate}

  \textbf{Ejemplo:} convertir $150 \text{ cm}$ a metros:
  $$150 \text{ cm} \cdot \frac{1 \text{ m}}{100 \text{ cm}} = \frac{150}{100} \text{ m} = 1{,}5 \text{ m}$$
\end{cuadropasos}

\begin{cuadroteoria}
  \textbf{Conversión compuesta km/h $\leftrightarrow$ m/s:} como $1 \text{ km} = 1000 \text{ m}$ y $1 \text{ h} = 3600 \text{ s}$:
  $$1 \frac{\text{km}}{\text{h}} = \frac{1000 \text{ m}}{3600 \text{ s}} = \frac{1}{3{,}6} \frac{\text{m}}{\text{s}}$$
  \begin{itemize}[leftmargin=*]
    \item Para pasar de \textbf{km/h a m/s}: se divide entre $3{,}6$.
    \item Para pasar de \textbf{m/s a km/h}: se multiplica por $3{,}6$.
  \end{itemize}
\end{cuadroteoria}

\begin{cuadroteoria}
  \textbf{Densidad:} es la masa contenida en cada unidad de volumen.
  $$\rho = \frac{m}{V} \qquad (\text{en } \text{kg/m}^3 \text{ o } \text{g/cm}^3)$$
  El agua tiene densidad $1000 \text{ kg/m}^3 = 1 \text{ g/cm}^3$.
\end{cuadroteoria}

\newpage %ajustes pag

\begin{cuadroadvertencia}
  \textbf{Errores clásicos:}
  \begin{itemize}[leftmargin=*]
    \item Multiplicar en lugar de dividir: de km/h a m/s se \emph{divide} entre $3{,}6$ (72 km/h $= 20$ m/s, no $259{,}2$).
    \item Cancelar unidades que no corresponden: el factor debe colocarse de modo que la unidad de partida quede \emph{abajo}.
    \item Olvidar las unidades en el resultado final: un número sin unidad no es una medida física.
  \end{itemize}
\end{cuadroadvertencia}

\subsection{Ejercicios resueltos}

\textbf{Ejemplo R1.} Convierte $150 \text{ cm}$ a metros.

\begin{cuadrosolucion}
  \textbf{Solución:}
  $$150 \text{ cm} \cdot \frac{1 \text{ m}}{100 \text{ cm}} = 1{,}5 \text{ m}.$$
\end{cuadrosolucion}

\textbf{Ejemplo R2.} Convierte $2{,}5$ horas a segundos.

\begin{cuadrosolucion}
  \textbf{Solución:} usamos dos factores (horas a minutos y minutos a segundos):
  $$2{,}5 \text{ h} \cdot \frac{60 \text{ min}}{1 \text{ h}} \cdot \frac{60 \text{ s}}{1 \text{ min}} = 2{,}5 \cdot 60 \cdot 60 \text{ s} = 9000 \text{ s}.$$
\end{cuadrosolucion}

\textbf{Ejemplo R3.} Un autobús va a $72$ km/h. ¿A cuántos m/s equivale?

\begin{cuadrosolucion}
  \textbf{Solución:} de km/h a m/s se divide entre $3{,}6$:
  $$72 \frac{\text{km}}{\text{h}} \div 3{,}6 = 20 \frac{\text{m}}{\text{s}}.$$
\end{cuadrosolucion}

\textbf{Ejemplo R4.} Un bloque de $2 \text{ kg}$ ocupa un volumen de $0{,}004 \text{ m}^3$. Calcula su densidad.

\begin{cuadrosolucion}
  \textbf{Solución:}
  $$\rho = \frac{m}{V} = \frac{2 \text{ kg}}{0{,}004 \text{ m}^3} = 500 \frac{\text{kg}}{\text{m}^3}.$$
\end{cuadrosolucion}

\newpage %ajuste para pág

\subsection{Ejercicios propuestos}

\begin{enumerate}[leftmargin=*, resume]
  \item \facil{} Convierte a metros: $a)\ 250 \text{ cm} \qquad b)\ 4 \text{ km} \qquad c)\ 500 \text{ mm}$.

  \item \facil{} Convierte a kilogramos: $a)\ 3500 \text{ g} \qquad b)\ 0{,}5 \text{ t}$ (1 t $= 1000$ kg).

  \item \facil{} Convierte a segundos: $a)\ 3 \text{ min} \qquad b)\ 1 \text{ h} \qquad c)\ 90 \text{ min}$.

  \item \facil{} Convierte a m/s: $a)\ 36 \text{ km/h} \qquad b)\ 90 \text{ km/h}$.

  \item \medio{} Convierte a km/h: $a)\ 10 \text{ m/s} \qquad b)\ 25 \text{ m/s}$.

  \item \medio{} Una maratón mide $42{,}195$ km. ¿Cuántos metros son? ¿Cuántos centímetros?

  \item \medio{} Un tren viaja a $144$ km/h. ¿Cuántos metros recorre en un segundo?

  \item \medio{} Calcula la densidad de un objeto de $400 \text{ g}$ que ocupa $50 \text{ cm}^3$. Exprésala en g/cm$^3$ y en kg/m$^3$.

  \item \dificil{} Un corredor avanza $100$ m en $9{,}58$ s (récord mundial). ¿Cuál es su rapidez media en m/s y en km/h?

  \item \dificil{} El agua tiene densidad $1 \text{ g/cm}^3$. Demuestra que equivale a $1000 \text{ kg/m}^3$ usando factores de conversión.

  \item \dificil{} Un coche recorre $240$ km en $2{,}5$ horas. ¿Cuál es su rapidez media en km/h y en m/s?
\end{enumerate}

\newpage

% ============================================================
\section{Vectores}
% ============================================================

\subsection{Escalares y vectores}

\begin{cuadroteoria}
  \begin{itemize}[leftmargin=*]
    \item \textbf{Magnitud escalar:} queda definida solo con un número y su unidad. Ejemplos: masa ($5$ kg), tiempo ($10$ s), temperatura ($30^{\circ}$C).
    \item \textbf{Magnitud vectorial:} necesita módulo, dirección y sentido. Ejemplos: desplazamiento, velocidad, fuerza.
  \end{itemize}

  Un \textbf{vector} se representa con una flecha y se denota $\vec{v}$:
  \begin{center}
    \begin{tikzpicture}[scale=1]
      \draw[thick, -Latex, mwprimario] (0,0) -- (3.5,1.75);
      \filldraw (0,0) circle (2pt) node[below left] {origen};
      \node[mwprimario] at (3.9,1.9) {$\vec{v}$};
      \draw[mwrojonaranja, thick, -Latex] (4.5,0) -- (4.5,2) node[right] {módulo (longitud)};
      \draw[gray, dashed] (0,0) -- (3.5,0);
      \draw[mwvioleta, thick] (0.7,0) arc (0:26.6:0.7);
      \node[mwvioleta] at (1.05,0.25) {$\theta$};
    \end{tikzpicture}
  \end{center}
  \begin{itemize}[leftmargin=*]
    \item \textbf{Módulo} $|\vec{v}|$: la longitud de la flecha (la "cantidad" del vector).
    \item \textbf{Dirección}: la recta sobre la que actúa (dada por el ángulo $\theta$).
    \item \textbf{Sentido}: hacia dónde apunta la flecha.
  \end{itemize}
\end{cuadroteoria}

\subsection{Componentes rectangulares}

\begin{cuadroteoria}
  Todo vector se puede descomponer en dos \textbf{componentes rectangulares} (una sobre el eje $X$ y otra sobre el eje $Y$):
  $$\vec{v} = (v_x, v_y) \qquad \text{con} \qquad v_x = |\vec{v}| \cos \theta, \quad v_y = |\vec{v}| \sen \theta$$
  \begin{center}
    \begin{tikzpicture}[scale=0.95]
      \draw[thin, gray!40] (-0.5,-0.5) grid (4.5,3.5);
      \draw[thick, -Latex] (-0.8,0) -- (4.8,0) node[right] {$x$};
      \draw[thick, -Latex] (0,-0.8) -- (0,3.8) node[above] {$y$};
      \draw[thick, -Latex, mwprimario] (0,0) -- (3.6,2.4) node[above right] {$\vec{v}$};
      \draw[mwprimario, thick] (0.7,0) arc (0:33.7:0.7);
      \node[mwprimario] at (1.0,0.25) {$\theta$};
      \draw[mwmagenta, thick, -Latex] (0,0) -- (3.6,0) node[midway, below] {$v_x$};
      \draw[mwmagenta, thick, -Latex] (3.6,0) -- (3.6,2.4) node[midway, right] {$v_y$};
      \draw[gray, dashed] (0,0) -- (0,2.4);
      \draw[gray, dashed] (3.6,2.4) -- (0,2.4);
    \end{tikzpicture}
  \end{center}
  Y al revés, conociendo las componentes se puede recuperar el módulo con Pitágoras:
  $$|\vec{v}| = \sqrt{v_x^2 + v_y^2}$$
\end{cuadroteoria}

\subsection{Suma de vectores}

\begin{cuadroteoria}
  \textbf{Método del polígono (gráfico):} se dibujan los vectores uno seguido de otro (el origen de cada uno en la punta del anterior). El vector resultante une el origen del primero con la punta del último.
  \begin{center}
    \begin{tikzpicture}[scale=0.9]
      \draw[thick, -Latex, mwprimario] (0,0) -- (2.5,1.2) node[midway, above] {$\vec{a}$};
      \draw[thick, -Latex, mwnaranja] (2.5,1.2) -- (4.5,0.4) node[midway, above] {$\vec{b}$};
      \draw[thick, -Latex, mwmagenta] (0,0) -- (4.5,0.4) node[midway, below] {$\vec{a} + \vec{b}$};
    \end{tikzpicture}
  \end{center}

  \textbf{Método analítico (por componentes):} se suman las componentes de cada eje por separado.
  $$\vec{a} = (a_x, a_y), \quad \vec{b} = (b_x, b_y) \quad \Rightarrow \quad \vec{a} + \vec{b} = (a_x + b_x, a_y + b_y)$$

  \textbf{Multiplicar un vector por un escalar $k$:} cambia el módulo (se alarga o acorta) y, si $k < 0$, invierte el sentido.
\end{cuadroteoria}

\begin{cuadropasos}
  \textbf{Sumar vectores por componentes:}
  \begin{enumerate}[leftmargin=*]
    \item Descompón cada vector en sus componentes con $\cos \theta$ y $\sen \theta$.
    \item Suma las componentes $x$ entre sí y las $y$ entre sí.
    \item Calcula el módulo de la resultante con Pitágoras.
  \end{enumerate}
\end{cuadropasos}

\begin{cuadroadvertencia}
  \textbf{Errores clásicos:}
  \begin{itemize}[leftmargin=*]
    \item Sumar vectores como si fueran números: $3$ km al este + $4$ km al norte NO son $7$ km: son $\sqrt{3^2 + 4^2} = 5$ km.
    \item Confundir dirección con sentido: "norte" es una dirección, "hacia el norte" es un sentido.
    \item Olvidar el signo de las componentes: un vector que apunta a la izquierda o hacia abajo tiene componentes negativas.
  \end{itemize}
\end{cuadroadvertencia}

\newpage %ajustes pag

\subsection{Reducción al primer cuadrante}

\begin{cuadroteoria}
  Cuando un vector forma un ángulo $\theta$ \emph{arbitrario} con el eje $X$ positivo, se habla de su \textbf{ángulo de referencia} $\theta_{\text{ref}}$ (siempre entre $0^{\circ}$ y $90^{\circ}$) y del \textbf{cuadrante} en que se ubica su punta.  El ángulo de referencia es el ángulo agudo formado con el eje $X$ (positivo o negativo, según el cuadrante).

  \begin{center}
    \begin{tikzpicture}[scale=1.05]
      % ejes
      \draw[thick, -Latex] (-3,0) -- (3,0) node[right] {$x$};
      \draw[thick, -Latex] (0,-3) -- (0,3) node[above] {$y$};
      % etiquetas de cuadrantes
      \node[mwprimario!80, font=\small\bfseries] at ( 1.6, 1.6) {I};
      \node[mwnaranja!90, font=\small\bfseries]  at (-1.6, 1.6) {II};
      \node[mwmagenta!90, font=\small\bfseries]  at (-1.6,-1.6) {III};
      \node[mwvioleta!90, font=\small\bfseries]  at ( 1.6,-1.6) {IV};
      % vectores ejemplo en cada cuadrante
      \draw[thick,-Latex,mwprimario]  (0,0) -- ( 1.8, 1.2) node[above right,font=\tiny] {$\vec{v}_{\text{I}}$};
      \draw[thick,-Latex,mwnaranja]   (0,0) -- (-1.8, 1.2) node[above left, font=\tiny] {$\vec{v}_{\text{II}}$};
      \draw[thick,-Latex,mwmagenta]   (0,0) -- (-1.8,-1.2) node[below left, font=\tiny] {$\vec{v}_{\text{III}}$};
      \draw[thick,-Latex,mwvioleta]   (0,0) -- ( 1.8,-1.2) node[below right,font=\tiny] {$\vec{v}_{\text{IV}}$};
    \end{tikzpicture}
  \end{center}

  \textbf{Signos de las componentes según el cuadrante:}
  \begin{center}
    \begin{tabular}{cccc}
      \toprule
      Cuadrante & $\theta$ (eje $+x$)              & $v_x$ & $v_y$ \\
      \midrule
      I         & $0^\circ < \theta < 90^\circ$    & $+$   & $+$   \\
      II        & $90^\circ < \theta < 180^\circ$  & $-$   & $+$   \\
      III       & $180^\circ < \theta < 270^\circ$ & $-$   & $-$   \\
      IV        & $270^\circ < \theta < 360^\circ$ & $+$   & $-$   \\
      \bottomrule
    \end{tabular}
  \end{center}

  Las fórmulas de las componentes \textbf{siempre se aplican con el ángulo de referencia $\theta_{\text{ref}}$}; el signo se asigna según el cuadrante:
  $$v_x = \pm |\vec{v}|\cos\theta_{\text{ref}}, \qquad v_y = \pm |\vec{v}|\sen\theta_{\text{ref}}$$
\end{cuadroteoria}

\newpage %ajuste pág

\begin{cuadropasos}
  \textbf{Cómo descomponer un vector en cualquier cuadrante:}
  \begin{enumerate}[leftmargin=*]
    \item Localiza el cuadrante en que está el vector a partir de $\theta$ (medido desde el eje $+x$ en sentido antihorario).
    \item Calcula el ángulo de referencia $\theta_{\text{ref}}$:
          \begin{itemize}
            \item Cuadrante I:   $\theta_{\text{ref}} = \theta$
            \item Cuadrante II:  $\theta_{\text{ref}} = 180^\circ - \theta$
            \item Cuadrante III: $\theta_{\text{ref}} = \theta - 180^\circ$
            \item Cuadrante IV:  $\theta_{\text{ref}} = 360^\circ - \theta$
          \end{itemize}
    \item Aplica $v_x = |\vec{v}|\cos\theta_{\text{ref}}$ y $v_y = |\vec{v}|\sen\theta_{\text{ref}}$.
    \item Asigna el signo correcto según el cuadrante (tabla de arriba).
  \end{enumerate}
\end{cuadropasos}

\begin{cuadroadvertencia}
  \textbf{Errores clásicos en la reducción al primer cuadrante:}
  \begin{itemize}[leftmargin=*]
    \item Usar $\theta$ directamente en $\cos$ y $\sen$ sin reducir: $\cos 150^\circ \neq \cos 150^\circ$ como positivo; hay que reconocer que da $-\cos 30^\circ$.
    \item Olvidar el signo de la componente: un vector en el III cuadrante tiene \emph{ambas} componentes negativas.
    \item Confundir el ángulo de referencia con el ángulo desde el eje $+y$: el ángulo de referencia siempre se mide desde el eje $X$ más cercano.
  \end{itemize}
\end{cuadroadvertencia}

\subsection{Ejercicios resueltos}

\textbf{Ejemplo R1.} Clasifica cada magnitud como escalar o vectorial: $a)$ $25^{\circ}$C \quad $b)$ $60$ km/h hacia el este \quad $c)$ $5$ kg \quad $d)$ fuerza de $10$ N hacia arriba.

\begin{cuadrosolucion}
  \textbf{Solución:}
  \begin{itemize}
    \item $a)$ Temperatura: escalar.
    \item $b)$ Velocidad con dirección: vectorial.
    \item $c)$ Masa: escalar.
    \item $d)$ Fuerza con dirección: vectorial.
  \end{itemize}
\end{cuadrosolucion}

\textbf{Ejemplo R2.} Un vector $\vec{v}$ tiene módulo $10$ y forma un ángulo de $60^{\circ}$ con el eje $X$ positivo. Halla sus componentes.

\begin{cuadrosolucion}
  \textbf{Solución:} usando los valores notables ($\cos 60^{\circ} = \tfrac{1}{2}$, $\sen 60^{\circ} = \tfrac{\sqrt{3}}{2}$):
  $$v_x = 10 \cos 60^{\circ} = 10 \cdot \frac{1}{2} = 5, \qquad v_y = 10 \sen 60^{\circ} = 10 \cdot \frac{\sqrt{3}}{2} = 5\sqrt{3}.$$
  Entonces $\vec{v} = (5,\, 5\sqrt{3})$.
\end{cuadrosolucion}

\textbf{Ejemplo R3.} Suma $\vec{a} = (3, 4)$ y $\vec{b} = (1, 2)$ y calcula el módulo de la resultante.

\begin{cuadrosolucion}
  \textbf{Solución:}
  $$\vec{a} + \vec{b} = (3 + 1, 4 + 2) = (4, 6), \qquad |\vec{a} + \vec{b}| = \sqrt{4^2 + 6^2} = \sqrt{16 + 36} = \sqrt{52} = 2\sqrt{13}.$$
\end{cuadrosolucion}

\textbf{Ejemplo R4.} Un vector $\vec{u}$ tiene módulo $10$ y forma un ángulo de $150^{\circ}$ con el eje $X$ positivo. Halla sus componentes.

\begin{cuadrosolucion}
  \textbf{Solución:} $150^{\circ}$ está en el \textbf{cuadrante II}, así que $v_x < 0$ y $v_y > 0$.
  \begin{enumerate}[leftmargin=*, label=\arabic*.]
    \item Ángulo de referencia: $\theta_{\text{ref}} = 180^{\circ} - 150^{\circ} = 30^{\circ}$.
    \item Componentes (con valores notables $\cos 30^{\circ} = \frac{\sqrt{3}}{2}$, $\sen 30^{\circ} = \frac{1}{2}$):
          $$v_x = -\,10\cos 30^{\circ} = -\,10 \cdot \frac{\sqrt{3}}{2} = -5\sqrt{3} \approx -8{,}66,$$
          $$v_y = +\,10\sen 30^{\circ} = +\,10 \cdot \frac{1}{2} = 5.$$
  \end{enumerate}
  Resultado: $\vec{u} = (-5\sqrt{3},\; 5)$.
\end{cuadrosolucion}

\textbf{Ejemplo R5.} Un vector $\vec{w}$ tiene módulo $8$ y apunta en la dirección $225^{\circ}$ (cuadrante III). Halla sus componentes.

\begin{cuadrosolucion}
  \textbf{Solución:} $225^{\circ}$ está en el \textbf{cuadrante III}: ambas componentes son negativas.
  \begin{enumerate}[leftmargin=*, label=\arabic*.]
    \item Ángulo de referencia: $\theta_{\text{ref}} = 225^{\circ} - 180^{\circ} = 45^{\circ}$.
    \item Con $\cos 45^{\circ} = \sen 45^{\circ} = \frac{\sqrt{2}}{2}$:
          $$w_x = -\,8\cos 45^{\circ} = -4\sqrt{2} \approx -5{,}66, \qquad w_y = -\,8\sen 45^{\circ} = -4\sqrt{2} \approx -5{,}66.$$
  \end{enumerate}
  Resultado: $\vec{w} = (-4\sqrt{2},\; -4\sqrt{2})$.
\end{cuadrosolucion}

\subsection{Ejercicios propuestos}

\begin{enumerate}[leftmargin=*, resume]
  \item \facil{} Clasifica como escalar o vectorial: $a)$ $10$ s \quad $b)$ desplazamiento de $20$ m al norte \quad $c)$ $9{,}8$ m/s$^2$ hacia abajo \quad $d)$ $120$ W.

  \item \facil{} Nombra los tres elementos de un vector y dibuja un vector de módulo $4$ en dirección horizontal hacia la derecha.

  \item \facil{} Un vector $\vec{v}$ tiene módulo $8$ y ángulo $0^{\circ}$. ¿Cuáles son sus componentes?

  \item \facil{} Suma: $\vec{a} = (2, 5)$ y $\vec{b} = (3, 1)$. Escribe la resultante como par ordenado.

  \item \facil{} Multiplica el vector $\vec{v} = (2, 3)$ por el escalar $k = 3$. ¿Qué vector obtienes?

  \item \medio{} Un vector tiene módulo $12$ y forma $30^{\circ}$ con el eje $X$. Halla sus componentes (usa $\cos 30^{\circ} = \tfrac{\sqrt{3}}{2}$ y $\sen 30^{\circ} = \tfrac{1}{2}$).

  \item \medio{} Un auto se desplaza $3$ km al este y luego $4$ km al norte. ¿Cuál es el módulo del desplazamiento total? ¿Es lo mismo que recorrer $7$ km? Explica.

  \item \medio{} Suma $\vec{a} = (5, -2)$ y $\vec{b} = (-3, 4)$ y calcula el módulo de la resultante.

  \item \medio{} El viento sopla con un vector $\vec{v} = (3, 4)$. ¿Cuál es su módulo? ¿Qué ángulo forma aproximadamente con el eje $X$?

  \item \dificil{} Un vector tiene componentes $v_x = 4$ y $v_y = 4$. ¿Cuál es su módulo y qué ángulo forma con el eje $X$?

  \item \dificil{} \textbf{El reto de los campeones:} un barco navega $30$ km hacia el este y luego $40$ km en dirección $60^{\circ}$ al norte del este. ¿A qué distancia del punto de partida queda?

        % ── Ejercicios de reducción al primer cuadrante ──────────
  \item \medio{} Halla las componentes de los siguientes vectores (módulo $20$):
        $a)$ $\theta = 120^{\circ}$ \quad $b)$ $\theta = 240^{\circ}$ \quad $c)$ $\theta = 315^{\circ}$.\\
        Indica en cada caso el cuadrante y el ángulo de referencia.

  \item \dificil{} Una fuerza de $50$ N actúa con un ángulo de $210^{\circ}$ respecto al eje $+x$. Halla sus componentes $F_x$ y $F_y$, y verifica que $\sqrt{F_x^2 + F_y^2} = 50$ N.
\end{enumerate}

\newpage

% ============================================================
\section{Clave de Respuestas}
% ============================================================

\subsection*{Sección 1: Sistema Internacional y medición (Ejercicios 1--11)}

\begin{enumerate}[leftmargin=*, label=\textbf{\arabic*.}]
  \setcounter{enumi}{0}
  \item a) Fundamental. \quad b) Derivada. \quad c) Fundamental. \quad d) Derivada.
  \item a) kg. \quad b) s. \quad c) m.
  \item $1000$ m; $0{,}01$ m; $0{,}001$ m.
  \item a) $3000$ m. \quad b) $7{,}5$ m. \quad c) $0{,}012$ m.
  \item a) $3 \times 10^{5}$. \quad b) $5 \times 10^{-4}$. \quad c) $2{,}5 \times 10^{7}$.
  \item a) $2 \times 10^{-7}$. \quad b) $9{,}46 \times 10^{12}$.
  \item $2500$ mm; $250$ cm.
  \item $3{,}4 \times 10^{2}$ m/s.
  \item a) 3. \quad b) 2 (los ceros a la izquierda no cuentan). \quad c) 4 (el cero final cuenta).
  \item $1 \times 10^{-10}$ m = $1$ \AA.
  \item $5\,970\,000\,000\,000\,000\,000\,000\,000$ kg; la Tierra cabe $\approx 333\,000$ veces en el Sol (en masa).
\end{enumerate}

\subsection*{Sección 2: Transformación de unidades (Ejercicios 12--22)}

\begin{enumerate}[leftmargin=*, label=\textbf{\arabic*.}]
  \setcounter{enumi}{11}
  \item a) $2{,}5$ m. \quad b) $4000$ m. \quad c) $0{,}5$ m.
  \item a) $3{,}5$ kg. \quad b) $500$ kg.
  \item a) $180$ s. \quad b) $3600$ s. \quad c) $5400$ s.
  \item a) $36 \div 3{,}6 = 10$ m/s. \quad b) $90 \div 3{,}6 = 25$ m/s.
  \item a) $10 \cdot 3{,}6 = 36$ km/h. \quad b) $25 \cdot 3{,}6 = 90$ km/h.
  \item $42\,195$ m; $4\,219\,500$ cm.
  \item $144 \div 3{,}6 = 40$ m/s, así que recorre $40$ m en 1 s.
  \item $\rho = \dfrac{400 \text{ g}}{50 \text{ cm}^3} = 8$ g/cm$^3$ $= 8000$ kg/m$^3$.
  \item $v \approx \dfrac{100}{9{,}58} \approx 10{,}4$ m/s $\approx 37{,}6$ km/h.
  \item $1 \frac{\text{g}}{\text{cm}^3} = \frac{10^{-3} \text{ kg}}{10^{-6} \text{ m}^3} = 10^{3} \frac{\text{kg}}{\text{m}^3} = 1000$ kg/m$^3$. \checkmark
  \item $v = \dfrac{240}{2{,}5} = 96$ km/h $= 96 \div 3{,}6 \approx 26{,}7$ m/s.
\end{enumerate}

\subsection*{Sección 3: Vectores (Ejercicios 23--35)}

\begin{enumerate}[leftmargin=*, label=\textbf{\arabic*.}]
  \setcounter{enumi}{22}
  \item a) Escalar. \quad b) Vectorial. \quad c) Vectorial. \quad d) Escalar.
  \item Módulo, dirección y sentido.
  \item $v_x = 8$, $v_y = 0$ (está sobre el eje $X$).
  \item $(5, 6)$.
  \item $(6, 9)$ (se triplican las componentes).
  \item $v_x = 12 \cdot \dfrac{\sqrt{3}}{2} = 6\sqrt{3}$; $v_y = 12 \cdot \dfrac{1}{2} = 6$. $\vec{v} = (6\sqrt{3}, 6)$.
  \item $|\vec{d}| = \sqrt{3^2 + 4^2} = 5$ km. No es lo mismo: la distancia recorrida es $7$ km, pero el desplazamiento es $5$ km.
  \item $(2, 2)$; módulo $\sqrt{8} = 2\sqrt{2}$.
  \item $|\vec{v}| = \sqrt{3^2 + 4^2} = 5$; el ángulo es aproximadamente $53^{\circ}$ (cerca de $60^{\circ}$, pero un poco menos).
  \item $|\vec{v}| = \sqrt{4^2 + 4^2} = 4\sqrt{2}$; ángulo $45^{\circ}$.
  \item \textbf{Reto:} segundo tramo: $(40\cos 60^{\circ},\, 40\sen 60^{\circ}) = (20,\, 20\sqrt{3})$. Total: $(30 + 20,\, 20\sqrt{3}) = (50,\, 20\sqrt{3})$. Distancia: $\sqrt{50^2 + (20\sqrt{3})^2} = \sqrt{2500 + 1200} = \sqrt{3700} \approx 60{,}8$ km.
  \item $a)$ C.II, $\theta_{\text{ref}} = 60^{\circ}$: $\vec{v} = (-10,\; 10\sqrt{3})$. \quad
        $b)$ C.III, $\theta_{\text{ref}} = 60^{\circ}$: $\vec{v} = (-10,\; -10\sqrt{3})$. \quad
        $c)$ C.IV, $\theta_{\text{ref}} = 45^{\circ}$: $\vec{v} = (10\sqrt{2},\; -10\sqrt{2})$.
  \item $\theta_{\text{ref}} = 210^{\circ} - 180^{\circ} = 30^{\circ}$ (C.III): $F_x = -50\cos 30^{\circ} = -25\sqrt{3} \approx -43{,}3$ N; $F_y = -50\sen 30^{\circ} = -25$ N. Verificación: $\sqrt{(-25\sqrt{3})^2 + (-25)^2} = \sqrt{1875 + 625} = \sqrt{2500} = 50$ N. \checkmark
\end{enumerate}

\vspace{8pt}
\begin{cuadroinfo}
  \small
  \textbf{¿Cómo te fue?} Ya sabes medir, convertir y trabajar con direcciones: eso es media física. En la próxima guía usaremos estas herramientas para describir el \emph{movimiento}: con el MRU y el MRUV aprenderás a predecir dónde estará un auto dentro de 10 segundos. ¡Nos vemos ahí!
\end{cuadroinfo}

\end{document}
